\documentclass{../base/canon}

%% canon_variables.tex — Valores por defecto de metadatos institucionales
%%
%% Incluir ANTES del \begin{document} para sobreescribir los defaults de la clase.
%% El template puede sobreescribir cualquier valor después de este \input.
%%
%% Uso en template:
%%   %% canon_variables.tex — Valores por defecto de metadatos institucionales
%%
%% Incluir ANTES del \begin{document} para sobreescribir los defaults de la clase.
%% El template puede sobreescribir cualquier valor después de este \input.
%%
%% Uso en template:
%%   %% canon_variables.tex — Valores por defecto de metadatos institucionales
%%
%% Incluir ANTES del \begin{document} para sobreescribir los defaults de la clase.
%% El template puede sobreescribir cualquier valor después de este \input.
%%
%% Uso en template:
%%   \input{../base/canon_variables}
%%   \SetDocTitle{Mi Informe Específico}   % sobreescribe el default

\SetDocLogo{../assets/logo.pdf}
\SetDocUnit{Tecnología \& Sistemas}
\SetContactUnit{Mesa de soporte}
\SetContactMembers{%
  Equipo CoreX & Soporte & soporte@empresa.com%
}
\SetContactNote{Para soporte técnico o consultas sobre este documento.}

%%   \SetDocTitle{Mi Informe Específico}   % sobreescribe el default

\SetDocLogo{../assets/logo.pdf}
\SetDocUnit{Tecnología \& Sistemas}
\SetContactUnit{Mesa de soporte}
\SetContactMembers{%
  Equipo CoreX & Soporte & soporte@empresa.com%
}
\SetContactNote{Para soporte técnico o consultas sobre este documento.}

%%   \SetDocTitle{Mi Informe Específico}   % sobreescribe el default

\SetDocLogo{../assets/logo.pdf}
\SetDocUnit{Tecnología \& Sistemas}
\SetContactUnit{Mesa de soporte}
\SetContactMembers{%
  Equipo CoreX & Soporte & soporte@empresa.com%
}
\SetContactNote{Para soporte técnico o consultas sobre este documento.}

\SetDocType{COMPONENTES}
\SetDocTitle{Catálogo Visual de Componentes Canon}
\SetDocCode{ANX-SYS-001}
\SetDocArea{Sistema Documental}
\SetDocAuthor{Equipo CoreX}
\SetDocDate{\today}
\SetDocUnit{Tecnología \& Sistemas}
\SetContactUnit{Soporte Documental}
\SetContactMembers{%
  Andrea Morales & Coordinadora & andrea.morales@empresa.com \\
  Luis Samaniego & Técnico PDF  & luis.samaniego@empresa.com%
}
\SetContactNote{Para consultas sobre implementación web del sistema canon.}

\begin{document}

%%═══════════════════════════════════════════════════════════════════════
\section{Jerarquía tipográfica}

\subsection{Subtítulo de sección (subsection)}

\subsubsection{Subsubtítulo de detalle (subsubsection)}

Texto de cuerpo normal con tono institucional: el desempeño operativo de la semana
reportó un nivel de cumplimiento superior al umbral establecido en el plan de trabajo,
con 97 de 100 lotes procesados dentro del tiempo objetivo.  Los indicadores de
productividad se mantuvieron estables y no se registraron incidencias críticas.

\medskip
\textbf{Negrita}: énfasis en términos clave y valores que requieren atención.

\medskip
\textit{Cursiva}: citas, términos técnicos o referencias documentales internas.

\medskip
\texttt{Monospace}: rutas de archivo, códigos de sistema, identificadores.
Ejemplo en línea: \CodePath{/api/balanzas/reporte?semana=17\&anio=2026}

%%═══════════════════════════════════════════════════════════════════════
\section{Ambientes de párrafo}

\subsection{Párrafo ejecutivo}

\begin{ParrafoEjecutivo}
Durante la semana 17-2026 se procesaron \textbf{842\,000 tallos} en el centro de
postcosecha, registrando un rendimiento de clasificación del \textbf{96.3\,\%} frente
al objetivo del 95\,\%.  El indicador de merma se situó en 2.1\,\%, dentro del rango
aceptable.  No se reportaron paros de línea no programados ni incidentes de seguridad.
\end{ParrafoEjecutivo}

\subsection{Párrafo metodológico}

\begin{ParrafoMetodologico}
Los datos presentados fueron extraídos del sistema ERP en el intervalo comprendido
entre el 2026-04-14 00:00:00 y el 2026-04-20 23:59:59 (UTC-5).  La agregación se
realizó por turno, área y variedad mediante la vista \texttt{v\_bal\_semana}.
Los totales parciales se reconciliarán con el sistema de nómina al cierre de ciclo.
\end{ParrafoMetodologico}

%%═══════════════════════════════════════════════════════════════════════
\section{Tablas}

\subsection{Tabla corta — tabular estándar}

\begin{table}[H]
  \centering
  \caption{Resumen de rendimiento por área — Semana 17-2026}
  \begin{tabular}{l r r r}
    \toprule
    \textbf{Área} & \textbf{Ingreso} & \textbf{Salida} & \textbf{Rend.\,\%} \\
    \midrule
    Apertura       & 280\,400 & 271\,300 & 96.8 \\
    Clasificación  & 271\,300 & 261\,200 & 96.3 \\
    Empaque        & 261\,200 & 257\,900 & 98.7 \\
    \midrule
    \textbf{Total} & \textbf{812\,900} & \textbf{790\,400} & \textbf{97.2} \\
    \bottomrule
  \end{tabular}
\end{table}

\subsection{Tabla ancha — adjustbox}

\begin{table}[H]
  \centering
  \caption{Detalle de mermas por variedad, turno y causa — Semana 17-2026}
  \begin{adjustbox}{max width=\linewidth}
  \begin{tabular}{l l l r r r r}
    \toprule
    \textbf{Variedad} & \textbf{Turno} & \textbf{Causa} &
    \textbf{Ingreso} & \textbf{Merma} & \textbf{Merma\,\%} & \textbf{Meta\,\%} \\
    \midrule
    Freedom   & Mañana & Tallo quebrado       & 48\,200 & 820  & 1.7 & 2.0 \\
    Freedom   & Tarde  & Botón abierto        & 41\,500 & 995  & 2.4 & 2.0 \\
    Explorer  & Mañana & Daño mecánico        & 35\,100 & 490  & 1.4 & 2.0 \\
    Explorer  & Tarde  & Longitud fuera rango & 31\,900 & 768  & 2.4 & 2.0 \\
    Vendela   & Mañana & Mancha botrytis      & 27\,300 & 628  & 2.3 & 2.0 \\
    Vendela   & Tarde  & Tallo quebrado       & 24\,800 & 397  & 1.6 & 2.0 \\
    Mondial   & Mañana & Botón abierto        & 22\,100 & 353  & 1.6 & 2.0 \\
    Mondial   & Tarde  & Longitud fuera rango & 19\,400 & 504  & 2.6 & 2.0 \\
    \midrule
    \multicolumn{3}{l}{\textbf{Total}} &
      \textbf{250\,300} & \textbf{4\,955} & \textbf{1.98} & 2.0 \\
    \bottomrule
  \end{tabular}
  \end{adjustbox}
\end{table}

\subsection{Tabla larga — longtable}

\begin{center}
\begin{longtable}{@{} l r r r r @{}}
  \caption{Registro de lotes procesados — Semana 17-2026} \\
  \toprule
  \textbf{Lote} & \textbf{Ingreso} & \textbf{Salida} & \textbf{Merma} & \textbf{Rend.\,\%} \\
  \midrule
  \endfirsthead
  \toprule
  \textbf{Lote} & \textbf{Ingreso} & \textbf{Salida} & \textbf{Merma} & \textbf{Rend.\,\%} \\
  \midrule
  \endhead
  \midrule
  \multicolumn{5}{r}{\small\itshape Continúa\ldots} \\
  \endfoot
  \bottomrule
  \endlastfoot
  LOT-2026-1401 & 12\,400 & 12\,016 & 384 & 96.9 \\
  LOT-2026-1402 & 11\,800 & 11\,446 & 354 & 97.0 \\
  LOT-2026-1403 & 13\,200 & 12\,804 & 396 & 97.0 \\
  LOT-2026-1404 & 10\,900 & 10\,573 & 327 & 97.0 \\
  LOT-2026-1405 & 14\,500 & 14\,065 & 435 & 97.0 \\
  LOT-2026-1406 & 11\,100 & 10\,767 & 333 & 97.0 \\
  LOT-2026-1407 & 12\,700 & 12\,319 & 381 & 97.0 \\
  LOT-2026-1408 &  9\,800 &  9\,506 & 294 & 97.0 \\
  LOT-2026-1409 & 15\,200 & 14\,744 & 456 & 97.0 \\
  LOT-2026-1410 & 11\,500 & 11\,155 & 345 & 97.0 \\
  LOT-2026-1411 & 10\,300 &  9\,991 & 309 & 97.0 \\
  LOT-2026-1412 & 13\,600 & 13\,192 & 408 & 97.0 \\
  LOT-2026-1413 & 12\,000 & 11\,640 & 360 & 97.0 \\
  LOT-2026-1414 & 14\,100 & 13\,677 & 423 & 97.0 \\
  LOT-2026-1415 & 11\,700 & 11\,349 & 351 & 97.0 \\
  \midrule
  \textbf{Total} & \textbf{184\,800} & \textbf{179\,244} & \textbf{5\,556} & \textbf{97.0} \\
\end{longtable}
\end{center}

%%═══════════════════════════════════════════════════════════════════════
\section{Figuras}

\subsection{Figura con ruta válida}

\FiguraConFallback[0.82\linewidth]{../assets/placeholder_chart.pdf}%
  {Distribución de tallos por variedad — Semana 17 (fuente: ERP CoreX)}

\subsection{Figura con ruta inválida (activa el placeholder automático)}

\FiguraConFallback[0.72\linewidth]{../assets/no_existe.pdf}%
  {Gráfico de mermas por turno — pendiente de generación}

\FiguraPlaceholder[0.60\linewidth]{Diagrama de flujo del proceso — en elaboración}

%%═══════════════════════════════════════════════════════════════════════
\section{Cajas canónicas}

\begin{ObservationBox}[Observación operativa]
  El turno tarde del día 2026-04-15 presentó un incremento del 18\,\% en merma
  por «botón abierto» respecto a la semana anterior.  Se asocia a la llegada de
  lotes de variedad Freedom con mayor madurez de campo.
\end{ObservationBox}

\begin{KeyFindingBox}[Hallazgo clave]
  El rendimiento de empaque del turno mañana superó el umbral de 98.5\,\% por primera
  vez en el ciclo 2026, alcanzando 98.7\,\%.  Confirma la efectividad del ajuste de
  calibración implantado el 2026-04-10.
\end{KeyFindingBox}

\begin{WarningBox}[Alerta de calidad]
  Se detectaron dos lotes (LOT-2026-1408, LOT-2026-1411) con merma superior al 3.0\,\%
  por causa «botrytis».  Requieren revisión de cadena de frío en origen.
\end{WarningBox}

\NoteInline{Nota:} Los datos de BAL2 del turno madrugada del 2026-04-17 no están
disponibles por mantenimiento. Se imputarán con promedio de los 3 días anteriores.

%%═══════════════════════════════════════════════════════════════════════
\section{Código}

\begin{CodeBlock}
-- Consulta de totales semanales
SELECT
    area,
    SUM(tallos_ingreso) AS ingreso,
    SUM(tallos_salida)  AS salida,
    ROUND(
        SUM(tallos_salida)::numeric
        / NULLIF(SUM(tallos_ingreso), 0) * 100, 1
    ) AS rendimiento_pct
FROM v_bal_semana
WHERE semana = 17 AND anio = 2026
GROUP BY area ORDER BY area;
\end{CodeBlock}

Endpoint: \CodePath{/api/balanzas/rendimiento?semana=17\&anio=2026}

%%═══════════════════════════════════════════════════════════════════════
\section{KPIs (FichaKPI)}

\begin{center}
\begin{tabular}{ccc}
  \FichaKPI{Tallos procesados}{842\,000}{semana 17-2026} &
  \FichaKPI{Rendimiento}{96.3\,\%}{meta: 95.0\,\%} &
  \FichaKPI{Merma total}{2.1\,\%}{límite: 3.0\,\%} \\[1.4em]
  \FichaKPI{Paros no prog.}{0}{sin incidentes} &
  \FichaKPI{Lotes}{15 / 15}{100\,\% completados} &
  \FichaKPI{Ciclo prom.}{4.2 min}{meta: ≤ 4.5 min} \\
\end{tabular}
\end{center}

%%═══════════════════════════════════════════════════════════════════════
\section{Listas}

\subsection{Hallazgos}

\begin{itemize}
  \item El 96.3\,\% de tallos cumplió los criterios de calidad del protocolo vigente.
  \item Las variedades de tallo corto presentaron mayor incidencia de «botón abierto».
  \item El tiempo promedio de ciclo en apertura fue 4.2 min (meta: ≤ 4.5 min).
  \item No se registraron paros ni incidentes de seguridad industrial.
\end{itemize}

\subsection{Recomendaciones}

\begin{enumerate}
  \item Revisar madurez de lotes Freedom antes del ingreso a clasificación.
  \item Evaluar ajuste del umbral de merma para variedades de tallo corto.
  \item Programar calibración preventiva del sistema BAL1 antes de semana 19.
  \item Documentar el procedimiento de imputación de datos faltantes.
\end{enumerate}

%%═══════════════════════════════════════════════════════════════════════
\section{Bloques de documento}

\subsection{MemoBlock}

\MemoBlock%
  {Gerencia General}%
  {Coordinación de Operaciones}%
  {Ejemplo de bloque memo para solicitudes y memorandos}%
  {MEM-OPS-000}

\subsection{AsistentesBlock y AcuerdosList}

\begin{AsistentesBlock}
  Asistente Uno  & Cargo A & asistente1@empresa.com \\
  Asistente Dos  & Cargo B & asistente2@empresa.com \\
\end{AsistentesBlock}

\begin{AcuerdosList}
  \item \textbf{Responsable — Fecha:} Descripción del primer acuerdo.
  \item \textbf{Responsable — Fecha:} Descripción del segundo acuerdo.
\end{AcuerdosList}

%%═══════════════════════════════════════════════════════════════════════
\section{Firma y contacto}

\SignatureBlock{Andrea Morales}{Coordinadora de Soporte Documental}

\sectionrule

\ContactSignature

\end{document}
